% Source control vertices transcribed verbatim from the user's TikZ message.
% This file preserves the original Hobby path; separate Python source
% uses periodic cubic splines and changes P10 y by +0.07 to regularize
% its coincidence with P3. Do not interpret this reference drawing as a
% certified KPC 3D model or transfer its crossing IDs without runtime audit.
\documentclass[tikz,border=6pt]{standalone}
% ------- TikZ Guide Lines -------
\newcommand{\SSTGuidesPoints}[2]{% #1=basename (e.g. P), #2=last index
  \ifsstguides
    \foreach \i in {1,...,#2}{
      \fill[blue] (#1\i) circle (1.2pt);
      \node[blue,font=\scriptsize,above] at (#1\i) {\i};
    }
    \draw[gray!40, dashed]
    \foreach \i [remember=\i as \lasti (initially 1)] in {2,...,#2,1} { (#1\lasti)--(#1\i) };
  \fi
}
% ------- TikZ Guide Lines -------

\usepackage{tikz}
\usetikzlibrary{hobby,knots}
\usetikzlibrary{decorations.markings}

\begin{document}
\begin{tikzpicture}[use Hobby shortcut]
\tikzset{knot diagram/every strand/.append style={draw=black}}
\coordinate (P1) at (0, 3.50);
\coordinate (P2) at (-3.25, 1.50);
\coordinate (P3) at (-3.25, 0);
\coordinate (P4) at (-1.50, -1.75);
\coordinate (P5) at (0.50, -1.25);
\coordinate (P6) at (-2, -0.50);
\coordinate (P7) at (-2.50, 1.50);
\coordinate (P8) at (0, 4);
\coordinate (P9) at (2.50, 1.50);
\coordinate (P10) at (2.25, -0.25);
\coordinate (P11) at (-0.50, -1.25);
\coordinate (P12) at (1.50, -1.75);
\coordinate (P13) at (3, -0.25);
\coordinate (P14) at (3.25, 1.50);

\begin{knot}[
    consider self intersections,
    clip width=5pt, clip radius=3pt,
    ignore endpoint intersections=false,
    flip crossing/.list={1,3,6,8},
    % ----draft mode=crossings % uncomment to see numbers
]
\strand
([closed] P1)..(P2)..(P3)..(P4)..(P5)..(P6)..(P7)..(P8)..(P9)..(P10)..(P11)..(P12)..(P13)..(P14);
\end{knot}
\SSTGuidesPoints{P}{14}
\end{tikzpicture}
\end{document}